\section*{第 \ref{ch:b2:limb-organogenesis} 章——器官发生：构建四足动物的肢}

\begin{solution}{exo:b2:limb-organogenesis:1}
近–远轴：沿肢端的\omterm{def:b2:limb-organogenesis:bud}{顶端外胚层嵴}，FGF。
前–后轴：后缘的\omterm{def:b2:limb-organogenesis:bud}{极性活性区}，
以 Sonic hedgehog 为信号。背–腹轴：背侧外胚层，Wnt7a。
\end{solution}

\begin{solution}{exo:b2:limb-organogenesis:2}
近段：肱骨、股骨（Hoxa/d 9--10）。中段：桡骨和尺骨，
胫骨和腓骨（Hox 11）。远段：腕和手，踝和足
（Hox 13）。
\end{solution}

\begin{solution}{exo:b2:limb-organogenesis:3}
肢区已确定（Tbx5）；\omterm{def:b2:limb-organogenesis:bud}{肢芽}在\omterm{def:b2:limb-organogenesis:bud}{顶端外胚层嵴}的 FGF 作用下伸长；
ZPA 和背侧外胚层使其形成模式；离开前进区
的细胞凝聚为指条；软骨分化；关节形成于
软骨被抑制之处；血管侵入，骨取代软骨；
指间细胞死亡；来自\omterm{prop:b2:vertebrate-development:neurulation}{体节}的肌肉前体细胞
迁入、融合，并借肌腱附着。
\end{solution}

\begin{solution}{exo:b2:limb-organogenesis:4}
（a）仅有很少或没有肱骨的残肢；（b）有肱骨、桡骨和尺骨，
但无手；（c）完整的肢——FGF 替代了\omterm{def:b2:limb-organogenesis:bud}{顶端外胚层嵴}。
\end{solution}

\begin{solution}{exo:b2:limb-organogenesis:5}
$x = 40\ln 2 = \qty{28}{\micro m}$，$40\ln 5 = \qty{64}{\micro m}$，
$40\ln 20 = \qty{120}{\micro m}$：区域宽度依次为 28、36 和
\qty{56}{\micro m}。
\end{solution}

\begin{solution}{exo:b2:limb-organogenesis:6}
（a）4-3-2-2-3-4，即完整的镜像组。（b）较弱的源产生
较低的峰值，因此仅增加低阈值的指：2-2-3-4。（c）
同样，短暂暴露仅确定最前端的额外
指：2-2-3-4——细胞随时间对浓度进行整合。
\end{solution}

\begin{solution}{exo:b2:limb-organogenesis:7}
$5000\times 2^{8} = 1.3\times 10^{6}$。达到 $4\times 10^{6}$ 需要
$\log_2 800 = 9.6$ 次倍增，因此细胞周期平均必须约为
\qty{10}{h}，或计数必须包括
从\omterm{prop:b2:vertebrate-development:neurulation}{体节}迁入的肌肉前体细胞；细胞死亡只会
使差异更大，而非减小。
\end{solution}

\begin{solution}{exo:b2:limb-organogenesis:8}
鸡的指间细胞接受 BMP 信号并死亡；鸭的
指间细胞表达 BMP 抑制剂并存活，因此蹼
得以保留。珠上的抑制剂挽救鸡蹼，使之成为
蹼足。移植体中按时死亡说明细胞在移除前已
接受指令——决定在早期作出，
并在之后自主执行。
\end{solution}

\begin{solution}{exo:b2:limb-organogenesis:9}
缺失 Hoxa13 和 Hoxd13 时：无远段——前肢止于
腕部。缺失 Hoxa11 和 Hoxd11 时：桡骨和尺骨明显缩短，
手几乎直接连于肱骨。每一段
都需要自己的一对 Hox 基因；这些结构域不仅仅是标记，
而是指令。
\end{solution}

\begin{solution}{exo:b2:limb-organogenesis:10}
来自嵴的 FGF 使 ZPA 中的 Shh 保持开启；Shh 经由中胚层
接力，使嵴中的 FGF 保持开启。因此每个中心的信号
可证明另一中心在起作用，因此肢体绝不会在未被
模式化时生长，或在未生长时被模式化，
任一方的暂时缺失可由另一方修复。若 Shh 不能维持
FGF，嵴将在一天内退化，生长将停于
已达到的任一节段：形成截短的肢。分娩的
正反馈循环由其驱动的事件终止——分娩移除了
刺激；肢体的循环则由分化终止——
中胚层在肢体完成时停止传递信号。
\end{solution}

\begin{solution}{exo:b2:limb-organogenesis:11}
$e^{4r} = 8$：$r = \ln 8/4 = \qty{0.52}{d^{-1}}$。生长：
指软骨增殖更快、时间更长（一种肢体增强子，可
提高指中生长基因的表达）。死亡：
指间组织通过表达 BMP 抑制剂得以存活，
从而保留膜。
\end{solution}

\begin{solution}{exo:b2:limb-organogenesis:12}
信号被重复使用，因为一个有效的受体–\omterm{def:b2:genome-diversification:hgt}{转导}模块
重新部署的代价低：在组织之间变化的不是
信号本身，而是其接收细胞有能力
开启的基因集合，因此相同的浓度分布在一处
意味着运动神经元，在另一处则意味着第 4 指。此类基因的\omterm{def:b2:genome-diversification:mutation}{突变}
因此会同时影响许多器官——Shh \omterm{def:b2:genome-diversification:mutation}{突变}同时造成
全前脑无裂畸形与肢体缺陷——除非\omterm{def:b2:genome-diversification:mutation}{突变}位于
组织特异性增强子中，那样它只改变一个器官，
这正是这些基因中多数演化变化的发生方式。
\end{solution}

\begin{solution}{pb:b2:limb-organogenesis:1}
\textbf{1.} $96/12 = 8$ 次倍增：$2000\times 256 = 5.1\times 10^{5}$
个细胞。
\textbf{2.} $1000/256 \approx 4$：细胞增大和软骨
基质。
\textbf{3.} 长度 $\times 10$；半径不变的圆柱体
体积也增至 $\times 10$；额外的百倍说明半径
也增长了约十倍——芽在伸长的同时
变宽成桨状。
\textbf{4.} 第 3 天时 $300/400 = \qty{75}{\%}$；第 7 天时 $300/4000 =
\qty{7.5}{\%}$。
\textbf{5.} 只有末端处于信号作用之下；随着末端推进，
留在后方的细胞按产生顺序离开该区域，因此
近端结构先被确定，远端结构最后被确定。
\textbf{6.} 中段（桡骨和尺骨）确定于第 3.5 至
4 天之间。
\textbf{7.} 近段于第 3.5 天前确定，中段于第 4.0 天前确定，远段于第 4.5
天前确定（指尖稍晚）。
\textbf{8.} 半天：一个细胞周期。
\textbf{9.} 每在前进区中经历一个周期，Hox 编码便推进一级
（9--10，然后 11，然后 13）；经过一、二或
三个周期后离开的细胞携带近段、中段或远段的编码。
\textbf{10.} 完整的翼：FGF 正是嵴所提供的。
\textbf{11.} 背侧表面长出的第二个芽：远端结构
重复，形成多余的肢。
\textbf{12.} 中胚层停止向嵴传递 Shh，Shh
表达终止，嵴退化，细胞
转而分化而非分裂：维持生长的循环被打破。
\textbf{13.} $k = \ln 2/1740 = \qty{4.0e-4}{s^{-1}}$；$\lambda =
\sqrt{10^{-12}/4\times 10^{-4}} = \qty{50}{\micro m}$。
\textbf{14.} $x = 50\ln(1/0.6) = 26$，$50\ln 4 = 69$，$50\ln 12.5 =
\qty{126}{\micro m}$：区域宽度分别为 26、43 和 \qty{57}{\micro m}。
\textbf{15.} $600 - 126 = \qty{474}{\micro m}$。
\textbf{16.} 每个边界向外移动 $50\ln 2 = \qty{35}{\micro
m}$：61、104、161。三个宽度分别为 61、43 和 \qty{57}{\micro m}：
只有第 4 指的区域增大，其他区域仅仅是被移位。
模式向前移动，而非增加一指，
无指的前部区域从 474 缩小到 \qty{439}{\micro m}。
\textbf{17.} 每组需要 \qty{126}{\micro m}：至少
\qty{252}{\micro m}；\qty{600}{\micro m} 的芽有空间，
中间有 \qty{348}{\micro m} 的无指组织。
\textbf{18.} 在 \qty{200}{\micro m} 处，两个梯度重叠并
相加：中间细胞接收到的信号足以形成第 3 指，第 2 指的
区域消失——4-3-4 或 4-3-3-4，指数少于完整的
镜像组。
\textbf{19.} 四分之一的细胞给出 $c_0/4$：没有细胞达到
第 4 指或第 3 指的阈值（边缘处即为 $0.25 c_0$），并且
第 2 指的阈值则外推至 $50\ln(0.25/0.08) =
\qty{57}{\micro m}$：多出一个第 2 指。
\textbf{20.} $3\times 300\times 300\times 200 = 5.4\times 10^{7}$
\unit{\micro m^3}：$5.4\times 10^{4}$ 个细胞；在 \qty{86400}{s} 内，
平均每秒 $0.6$ 个细胞。
\textbf{21.} 鸭蹼中的 BMP 珠：蹼死亡，趾
分开；鸡蹼中的抑制剂使蹼存活，形成
蹼状足。
\textbf{22.} $e^{0.5\times 4} = e^{2} = 7.4$。第二个变化是
BMP 抑制剂使指间膜存活。
\textbf{23.} 决定远段的 Hox 13 表达后期，以及
嵴的持续存在：鱼的褶皱形成
鳍条而非肢。\emph{Tiktaalik} 的鳍含有类腕骨——
一个正开始向肢转变的鳍。
\textbf{24.} $5.4\times 10^{4}/1000 = 54$：$\log_2 54 = 5.8$ 次倍增，
约 \qty{69}{h}，生成一天所去除的量需要三天。
\textbf{25.} 八次倍增；指的边界位于 26、69 和
\qty{126}{\micro m}；镜像重复至少需要
\qty{252}{\micro m}。
\end{solution}
